\chapter{Applicazioni web con Node.js}

Node.js ci permette di realizzare in modo semplice ed efficiente un server HTTP,
sfruttando i moduli integrati e il suo event loop a thread singolo. Node.js fa
anche un po' di pre-elaborazione sulle richieste (analizza gli header) e ci
fornisce un'astrazione per richieste e risposte.

A parte questo, però, siamo da soli. Per costruire una vera applicazione
dobbiamo affrontare tre problemi:

\begin{itemize}
  \item il \term{routing};
  \item il \term{templating};
  \item l'\term{analisi del corpo} delle richieste.
\end{itemize}

Li affrontiamo uno alla volta, costruendo la stessa applicazione
\guillemotleft lista di cose da fare\guillemotright\ già vista nell'esempio CGI.

\section{L'applicazione da costruire}

L'applicazione ha tre pagine. Vedremo se costruirla con Node.js sarà meno
doloroso del famigerato esempio con CGI e Bash.

\begin{itemize}
  \item La \textbf{homepage} è la pagina d'ingresso, con i collegamenti alla
        lista e alla pagina di azzeramento.
  \item La \textbf{lista} contiene un form per salvare nuovi elementi, inviato
        con \code{POST} \textbf{allo stesso URL} della pagina, e l'elenco degli
        elementi già salvati.
  \item La pagina di \textbf{azzeramento} (\textit{reset}) cancella tutti gli
        elementi salvati.
\end{itemize}

\noindent\begin{minipage}[t]{0.31\linewidth}
\begin{browser}[localhost:3000/]
  {\Lora\bfseries\large To-do list}\\[3pt]
  Benvenuto! Che cosa vuoi fare?\\[3pt]
  \textbullet\ \wtlink{Vedere la lista}\\
  \textbullet\ \wtlink{Azzerare l'app}
\end{browser}
\end{minipage}\hfill
\begin{minipage}[t]{0.35\linewidth}
\begin{browser}[localhost:3000/todo]
  {\Lora\bfseries\large La lista}\\[3pt]
  \wtcampo[1.9cm]{Learn Node.js}\ \wtbottone{Salva}\\[3pt]
  \textbullet\ Learn HTML\\
  \textbullet\ Learn CSS
\end{browser}
\end{minipage}\hfill
\begin{minipage}[t]{0.30\linewidth}
\begin{browser}[localhost:3000/reset]
  {\Lora\bfseries\large Azzerata}\\[3pt]
  Tutti gli elementi sono stati cancellati.\\[3pt]
  \wtlink{Torna alla lista}
\end{browser}
\end{minipage}

\section{Lo scheletro dell'applicazione}

\begin{affianco}[0.5][c]
\begin{lstlisting}[style=js, name=app.js]
import http from 'http';

const PORT = 3000;

let server = http.createServer();
server.listen(PORT);
server.on('request', handleRequest);

function handleRequest(request, response) {
  /* gestione della richiesta */
}
\end{lstlisting}
\risultato
\begin{immagine}[A ogni richiesta il server emette l'evento\\
                 \code{request}, gestito da \code{handleRequest}.]
\begin{tikzpicture}[
    a/.style={-{Stealth[length=5pt,width=4pt]}, line width=0.8pt, draw=neutral!45!black}]
  \node[mnota] (br) at (0,0) {richiesta HTTP};
  \node[wtnode, font=\FiraCodeNL\scriptsize] (sv) at (2.6,0) {server};
  \node[mfun] (h) at (6.3,0) {handleRequest(req, res)};
  \draw[a] (br) -- (sv);
  \draw[a, draw=accent!70!black] (sv) -- node[mcod, above=2pt]{'request'} (h);
\end{tikzpicture}
\end{immagine}
\end{affianco}

Rispetto all'esempio del capitolo precedente si nota una differenza stilistica:
invece di passare la callback direttamente a \code{createServer}, la si registra
come gestore dell'evento \code{request}. È lo stesso meccanismo degli eventi
del browser visto nel capitolo 7.

\section{Routing}

La nostra applicazione dovrà gestire molti tipi diversi di richiesta HTTP:
mostrare la homepage, mostrare la lista, azzerare l'applicazione, e così via.

\dfn{Routing}{
  Il \term{routing} è il problema di stabilire, data una richiesta, \textbf{quale
  codice} debba essere eseguito per servirla.}

\begin{figure}[H]
  \centering
  \begin{tikzpicture}[font=\Lato\small,
      n/.style={wtnode, minimum width=3.0cm, minimum height=1.0cm},
      a/.style={-{Stealth[length=5pt,width=4pt]}, line width=0.85pt,
                draw=neutral!45!black},
      c/.style={font=\FiraCodeNL\scriptsize, anchor=east,
                text=neutral!25!black, align=right}]

    \node[c] (r1) at (-0.3,0.8)  {GET /todos HTTP/1.1};
    \node[c] (r2) at (-0.3,-0.8) {POST /todos HTTP/1.1\\ todo=Learn+Node.js};

    \node[n, draw=orange-gradient-6, fill=orange-gradient-1!40!white]
      (rt) at (2.9,0) {\textbf{Router}};

    \node[wtnodeplain, anchor=west, font=\FiraCodeNL\scriptsize]
      (h1) at (5.3,0.8) {handleGetReqs()};
    \node[wtnodeplain, anchor=west, font=\FiraCodeNL\scriptsize]
      (h2) at (5.3,-0.8) {handlePostReqs()};

    \draw[a] (r1.east) -- (rt.west |- r1);
    \draw[a] (r2.east) -- (rt.west |- r2);
    \draw[a] (rt.east |- h1) -- (h1.west);
    \draw[a] (rt.east |- h2) -- (h2.west);
  \end{tikzpicture}
  \caption{Il router smista ogni richiesta verso il codice che la serve.}
  \label{fig:router}
\end{figure}

\subsection{Instradare in base al percorso}

\begin{affianco}[0.5][c]
\begin{lstlisting}[style=js, name=router.js]
function handleRequest(request, response) {
  switch (request.url) {
    case "/":
      renderHomepage(request, response);
      break;
    case "/reset":
      handleResetRequest(request, response);
      break;
    case "/todo":
      handleTodoListRequest(request, response);
      break;
    /* e cosi' via... */
    default:
      handleError(request, response, 404);
  }
}
\end{lstlisting}
\risultato
\begin{immagine}[La tabella di instradamento realizzata dallo \code{switch}.]
\begin{tikzpicture}[
    u/.style={mcod, anchor=east, draw=primary!55!black, fill=primary!8!white,
      rounded corners=2pt, minimum width=1.5cm, inner sep=3pt},
    f/.style={mcod, anchor=west},
    a/.style={-{Stealth[length=4pt,width=3.5pt]}, line width=0.7pt, draw=neutral!45!black}]
  \foreach \u/\f/\c [count=\i] in {\slash/renderHomepage/black,
      \slash reset/handleResetRequest/black,
      \slash todo/handleTodoListRequest/black, altro/handleError 404/merr} {
    \node[u] (u\i) at (1.5,-\i*0.6) {\u};
    \node[f, text=\c] (f\i) at (2.2,-\i*0.6) {\f};
    \draw[a] (u\i.east) -- (f\i.west);
  }
\end{tikzpicture}
\end{immagine}
\end{affianco}

\subsection{Instradare anche in base al metodo}

Le richieste verso un certo percorso possono essere gestite in modo diverso a
seconda del \textbf{metodo HTTP} usato:

\begin{itemize}
  \item una \code{GET} su \code{/todos} restituisce la lista degli elementi;
  \item una \code{POST} su \code{/todos} prova a salvarne uno nuovo.
\end{itemize}

\begin{affianco}[0.56][c]
\codicepiccolo
\begin{lstlisting}[style=js, name=router.js]
function handleTodoListRequest(request, response, context = {}) {
  switch (request.method) {
    case "GET":
      handleTodoListRequestGet(request, response, context); break;
    case "POST":
      handleTodoListRequestPost(request, response, context); break;
    default:
      handleError(request, response, 405, "Metodo non supportato");
  }
}
\end{lstlisting}
\risultato
\begin{immagine}[Stesso percorso, codice diverso a seconda del metodo.]
\begin{tikzpicture}[
    u/.style={mcod, draw=primary!55!black, fill=primary!8!white,
      rounded corners=2pt, inner sep=3pt},
    m/.style={mcod, font=\FiraCodeNL\tiny\bfseries, text=accent!70!black},
    f/.style={mcod, anchor=west},
    a/.style={-{Stealth[length=4pt,width=3.5pt]}, line width=0.7pt, draw=neutral!45!black}]
  \node[u] (t) at (0,0) {/todo};
  \foreach \m/\f/\y/\c in {GET/\ldots Get()/0.7/black, POST/\ldots Post()/0/black,
      altro/405/-0.7/merr} {
    \node[f, text=\c] (x) at (2.1,\y) {\f};
    \draw[a] (t.east) -- node[m, above, sloped]{\m} (x.west);
  }
\end{tikzpicture}
\end{immagine}
\end{affianco}

\info{Il codice 405}{
  Quando il percorso esiste ma il metodo non è previsto, la risposta corretta
  non è un \code{404} (\textit{Not Found}) bensì un \code{405} (\textit{Method
  Not Allowed}). È una distinzione che i client e i \textit{proxy} sanno
  interpretare, e che torna utile soprattutto nelle API REST.}

\subsection{Servire i file statici}

Se il nostro output HTML include file statici (immagini, CSS, script), dobbiamo
configurare il router perché li serva.

\begin{affianco}[0.6][c]
\codicepiccolo
\begin{lstlisting}[style=js, name=router.js]
switch (request.url) {
  /* ... */
  case "/css/bootstrap.css":
    fs.readFile('./static/css/bootstrap.css', function (err, data) {
      response.writeHead(200, { 'Content-Type': 'text/css' });
      response.end(data, 'utf-8');
    });
    break;

  case "/img/node.png":
    fs.readFile('./static/img/node.png', function (err, data) {
      response.writeHead(200, { 'Content-Type': 'image/png' });
      response.end(data, "binary");
    });
    break;
}
\end{lstlisting}
\risultato
\begin{immagine}[Ogni URL statica corrisponde a un file sul disco\\
                 e a un \code{Content-Type}.]
\begin{tikzpicture}[
    u/.style={mcod, anchor=west, draw=primary!55!black, fill=primary!8!white,
      rounded corners=2pt, inner sep=3pt},
    a/.style={-{Stealth[length=4pt,width=3.5pt]}, line width=0.7pt, draw=neutral!45!black}]
  \node[u] (u1) at (0,0) {/css/bootstrap.css};
  \node[mcod, anchor=west] (f1) at (0.1,-0.55) {./static/css/\ldots};
  \node[mnota, anchor=west] at (0.1,-0.95) {\code{text/css}};
  \node[u] (u2) at (0,-1.6) {/img/node.png};
  \node[mcod, anchor=west] (f2) at (0.1,-2.15) {./static/img/\ldots};
  \node[mnota, anchor=west] at (0.1,-2.55) {\code{image/png}};
\end{tikzpicture}
\end{immagine}
\end{affianco}

Si noti che \code{fs.readFile} è \textbf{asincrona}: prende una callback che
verrà invocata quando il file sarà stato letto. È esattamente il modello non
bloccante descritto nel capitolo precedente.

\subsection{Gestire gli errori}

Quando nessuna delle rotte note si applica, restituiamo un \code{404}.

\begin{affianco}[0.58][c]
\codicepiccolo
\begin{lstlisting}[style=js, name=router.js]
switch (request.url) {
  /* altre rotte */
  default:
    handleError(request, response, 404, "Pagina non trovata");
}

function handleError(request, response, statusCode, message) {
  let renderedContent = pug.renderFile("./templates/errorPage.pug",
    { "code": statusCode, "description": message });
  response.writeHead(statusCode, { "Content-Type": "text/html" });
  response.end(renderedContent);
}
\end{lstlisting}
\risultato
\begin{browser}[localhost:3000/nonesiste]
  {\Lora\bfseries\Huge 404}\\[4pt]
  Pagina non trovata
\end{browser}
\wtnota{La pagina d'errore, generata\\ dal template \code{errorPage.pug}.}
\end{affianco}

\warning{Il routing non resta semplice a lungo}{
  Uno \code{switch} sembra sufficiente, ma il problema cresce in fretta con
  l'applicazione:
  \begin{itemize}
    \item alcune rotte devono essere accessibili solo ad alcuni utenti (per
          esempio gli amministratori);
    \item alcune rotte dipendono da \textbf{schemi o parametri} nel percorso,
          come \code{GET /users/42/friends}, dove \code{42} può essere un id
          qualunque: con uno \code{switch} su una stringa esatta è impossibile;
    \item può servire eseguire \textbf{due o più pezzi di codice} per gestire una
          richiesta, per esempio una funzione che registra i dati e un'altra che
          prepara la risposta.
  \end{itemize}
  Sono esattamente i problemi che risolveremo con Express, nel prossimo
  capitolo.}

\section{Motori di template}

A un certo punto la nostra applicazione dovrà produrre del codice HTML da
rimandare nel corpo della risposta. Nel primo esempio l'abbiamo scritto a mano,
concatenando stringhe.

Il problema è che alcune parti dell'HTML si ripetono in più pagine (barra di
navigazione, piè di pagina) e che, man mano che le pagine diventano complesse,
costruire una risposta manipolando e concatenando stringhe diventa rapidamente
impraticabile.

\dfn{Motore di template}{
  Un \term{motore di template} (\textit{template engine}) permette di descrivere
  la struttura di una pagina in un file separato, con dei
  \textbf{segnaposto}, e di generare l'HTML finale combinando quel modello con i
  dati forniti a tempo di esecuzione.}

\subsection{Pug}

\term{Pug} (già noto come Jade) è un motore di template ad alte prestazioni
implementato in JavaScript, con una sintassi \textbf{sensibile agli spazi}
(simile a Python) per scrivere modelli HTML.

\begin{lstlisting}[style=shell]
npm install pug
\end{lstlisting}

\begin{lstlisting}[style=js]
import pug from "pug";

let html = pug.renderFile("./templates/basic_example.pug");
console.log(html);   // html contiene il codice HTML generato
\end{lstlisting}

\begin{esempio}{Il primo template Pug}
Il template a sinistra produce l'HTML a destra.

\begin{affianco}[0.45]
\begin{lstlisting}[style=pug, name=basic_example.pug]
doctype html
html(lang="it")
  head
    title Hello Pug!
  body
    h1(class="foo") Primo template Pug
    p Questo e' il nostro primo template Pug!
\end{lstlisting}
\risultato
\begin{lstlisting}[style=html, name=output.html]
<!DOCTYPE html>
<html lang="it">
  <head><title>Hello Pug!</title></head>
  <body>
    <h1 class="foo">Primo template Pug</h1>
    <p>Questo e' il nostro primo template Pug!</p>
  </body>
</html>
\end{lstlisting}
\end{affianco}

L'indentazione sostituisce i tag di chiusura, e gli attributi si scrivono fra
parentesi tonde. Il risultato è molto più breve e, soprattutto, impossibile da
sbagliare nell'annidamento.
\end{esempio}

\subsection{Inclusioni}

Pug è molto più di un modo diverso di scrivere HTML. Tanto per cominciare, un
template può \textbf{includerne} altri: è un grande aiuto per il riuso.

\begin{affianco}[0.55]
\begin{lstlisting}[style=pug, name=partials/head.pug]
head
  title Hello Pug
  script(src='/js/script.js')
  link(rel='stylesheet' href='/style.css')
\end{lstlisting}
\risultato
\begin{lstlisting}[style=pug, name=partials/footer.pug]
footer#footer
  p Copyright (c) Web Tech
\end{lstlisting}
\end{affianco}

\begin{affianco}[0.5][c]
\begin{lstlisting}[style=pug, name=index.pug]
doctype html
html
  include partials/head.pug
  body
    h1 Hello Pug!
    include partials/footer.pug
\end{lstlisting}
\risultato
\begin{immagine}[Gli \code{include} vengono sostituiti\\
                 dal contenuto dei file indicati.]
\begin{tikzpicture}[
    f/.style={draw=primary!55!black, fill=primary!6!white, rounded corners=2pt,
      line width=0.7pt, font=\FiraCodeNL\scriptsize, minimum width=2.6cm,
      minimum height=0.5cm},
    p/.style={f, draw=accent!55!black, fill=accent!12!white},
    a/.style={-{Stealth[length=4pt,width=3.5pt]}, line width=0.7pt, draw=accent!70!black}]
  \node[f] (r1) at (0,0) {html};
  \node[p] (r2) at (0,-0.6) {head (partials/head)};
  \node[f] (r3) at (0,-1.2) {body: h1 Hello Pug!};
  \node[p] (r4) at (0,-1.8) {footer (partials/footer)};
  \node[mcod, anchor=west] (h) at (2.4,-0.6) {head.pug};
  \node[mcod, anchor=west] (ft) at (2.4,-1.8) {footer.pug};
  \draw[a] (h.west) -- (r2.east);
  \draw[a] (ft.west) -- (r4.east);
\end{tikzpicture}
\end{immagine}
\end{affianco}

Il risultato è un'unica pagina, in cui i due frammenti hanno preso il posto
delle rispettive direttive \code{include}:

\begin{lstlisting}[style=html, name=output.html]
<!DOCTYPE html>
<html>
  <head>
    <title>Hello Pug</title>
    <script src="/js/script.js"></script>
    <link rel="stylesheet" href="/style.css"/>
  </head>
  <body>
    <h1>Hello Pug!</h1>
    <footer id="footer">
      <p>Copyright (c) Web Tech</p>
    </footer>
  </body>
</html>
\end{lstlisting}

Si noti \code{footer\#footer}: in Pug il cancelletto è l'abbreviazione per
l'attributo \attr{id}, esattamente come nei selettori CSS.

\subsection{Ereditarietà fra template}

Progettando applicazioni web è comune avere un \textbf{layout condiviso} da
tutte le pagine, che ciascuna pagina può poi personalizzare in alcune parti. Pug
supporta questo scenario con l'\term{ereditarietà fra template}, tramite le
parole chiave \code{block} ed \code{extends}.

\begin{figure}[H]
  \centering
  \begin{tikzpicture}[font=\Lato\scriptsize,
      b/.style={draw=primary!60!black, fill=primary!10!white, line width=0.7pt,
                minimum width=3.0cm, minimum height=0.5cm, anchor=north},
      bc/.style={b, draw=accent!60!black, fill=accent!12!white}]

    \foreach \x/\ttl in {0/{Layout base}, 4.4/{Homepage}, 8.8/{Pagina articolo}} {
      \node[font=\Lato\scriptsize\bfseries, text=neutral!25!black]
        at (\x,1.0) {\ttl};
    }

    % layout base
    \node[b] (h0) at (0,0.6) {heading};
    \node[b, below=1pt of h0] (n0) {navbar};
    \node[bc, below=1pt of n0, minimum height=1.1cm] (c0) {content};
    \node[b, below=1pt of c0] (f0) {footer};

    % homepage
    \node[b] (h1) at (4.4,0.6) {heading};
    \node[b, below=1pt of h1] (n1) {navbar};
    \node[bc, below=1pt of n1, minimum height=1.1cm] (c1) {contenuto proprio};
    \node[b, below=1pt of c1] (f1) {footer};

    % articolo
    \node[b] (h2) at (8.8,0.6) {heading};
    \node[b, below=1pt of h2] (n2) {navbar};
    \node[bc, below=1pt of n2, minimum height=1.1cm, minimum width=1.9cm,
          anchor=north west] (c2) at ([xshift=-1.5cm]n2.south west) {article};
    \node[bc, minimum height=1.1cm, minimum width=1.0cm, anchor=north west]
      (s2) at ([xshift=1pt]c2.north east) {sidebar};
    \node[b, below=1pt of c2.south west, anchor=north west] (f2) {footer};
  \end{tikzpicture}
  \caption{Un layout comune, personalizzato dalle singole pagine.}
  \label{fig:pug-inheritance}
\end{figure}

\begin{esempio}{Layout base e pagina figlia}
A sinistra il layout: definisce tre \textbf{blocchi}, \code{heading},
\code{content} e \code{footer}, e i blocchi \code{heading} e \code{footer}
hanno un contenuto predefinito. A destra la pagina figlia, che lo estende.

\begin{affianco}[0.5]
\begin{lstlisting}[style=pug, name=base-layout.pug]
doctype html
html
  head
    title Pug Inheritance
  body
    block heading
      h1 Pug Inheritance
    block content
    block footer
      footer Questo e' il footer predefinito.
\end{lstlisting}
\risultato
\begin{lstlisting}[style=pug, name=index.pug]
extends ./base-layout.pug

block content
  p Il contenuto vero e proprio della homepage.

block footer
  footer Questo e' un footer specializzato.
\end{lstlisting}
\end{affianco}

Il template figlio \textbf{sovrascrive} i blocchi \code{content} e
\code{footer}; il blocco \code{heading} non viene sovrascritto e mostra il
contenuto predefinito. Il risultato è:

\begin{affianco}[0.6][c]
\begin{lstlisting}[style=html, name=output.html]
<!DOCTYPE html>
<html>
  <head><title>Pug Inheritance</title></head>
  <body>
    <h1>Pug Inheritance</h1>
    <p>Il contenuto vero e proprio della homepage.</p>
    <footer>Questo e' un footer specializzato.</footer>
  </body>
</html>
\end{lstlisting}
\risultato
\begin{browser}[Pug Inheritance]
  {\Lora\bfseries\LARGE Pug Inheritance}\\[6pt]
  Il contenuto vero e proprio della homepage.\\[6pt]
  Questo e' un footer specializzato.
\end{browser}
\end{affianco}
\end{esempio}

\subsection{Locals: passare i dati al template}

I template Pug possono generare contenuto a partire da un oggetto di dati
fornito, chiamato \term{locals}, passato alle funzioni \code{render()} o
\code{renderFile()}.

\begin{esempio}{Interpolazione}
\begin{lstlisting}[style=js]
import pug from "pug";

let html = pug.renderFile("./locals.pug", {
  "product": {
    "name": "Samsung S24",
    "description": "Smartphone"
  }
});

console.log(html);
\end{lstlisting}

\begin{affianco}[0.55]
\begin{lstlisting}[style=pug, name=locals.pug]
- let footer = "una bella stringa"

h1 #{product.name.toUpperCase()}
p Descrizione: #{product.description}
footer #{footer}
\end{lstlisting}
\risultato
\begin{lstlisting}[style=html, name=output.html]
<h1>SAMSUNG S24</h1>
<p>Descrizione: Smartphone</p>
<footer>una bella stringa</footer>
\end{lstlisting}
\end{affianco}

Il codice contenuto fra \code{\#\{} e \code{\}} viene valutato,
\textbf{sottoposto a escaping} e inserito nell'output. L'escaping automatico
non è un dettaglio: è ciò che impedisce che un nome di prodotto contenente
\code{<script>} finisca per essere eseguito dal browser.

Una riga che inizia con \code{-} contiene invece codice JavaScript eseguito ma
non stampato.
\end{esempio}

\subsection{Condizionali}

\begin{lstlisting}[style=js]
let user = { "isAdmin": true, "name": "Brendan" };
let html = pug.renderFile("./conditionals.pug", { "user": user });
\end{lstlisting}

\begin{affianco}[0.55][c]
\begin{lstlisting}[style=pug, name=conditionals.pug]
if !user
  a(href="/login") Per favore, autenticati.
else if user.isAdmin
  p Ai tuoi ordini, #{user.name}
else
  p Bentornato, #{user.name}
\end{lstlisting}
\risultato
\begin{browser}[con user.isAdmin = true]
  Ai tuoi ordini, Brendan
\end{browser}
\wtnota{Solo il ramo scelto finisce nell'HTML.}
\end{affianco}

\subsection{Iterazioni}

\begin{lstlisting}[style=js]
let items = [
  { "name": "Margherita",  "price": 5.50 },
  { "name": "Marinara",    "price": 5.00 },
  { "name": "Capricciosa", "price": 6.50 }
];
let html = pug.renderFile("./iteration.pug", { "items": items });
\end{lstlisting}

\begin{affianco}[0.55][c]
\begin{lstlisting}[style=pug, name=iteration.pug]
ul
  each item in items
    li #{item.name} - EUR #{item.price.toFixed(2)}
  else
    li Nessuna pizza disponibile al momento!
\end{lstlisting}
\risultato
\begin{browser}[iteration]
  \begin{itemize}[leftmargin=1.8em, itemsep=0pt, topsep=2pt, parsep=0pt]
    \item Margherita - EUR 5.50
    \item Marinara - EUR 5.00
    \item Capricciosa - EUR 6.50
  \end{itemize}
\end{browser}
\end{affianco}

\begin{lstlisting}[style=html, name=output.html]
<ul>
  <li>Margherita - EUR 5.50</li>
  <li>Marinara - EUR 5.00</li>
  <li>Capricciosa - EUR 6.50</li>
</ul>
\end{lstlisting}

Il ramo \code{else} di un \code{each} viene eseguito quando la collezione è
vuota: è una comodità che evita di scrivere a mano il controllo.

\info{Non solo Pug}{
  Pug è solo uno dei tanti motori di template disponibili. Altri molto noti sono
  \term{EJS} (con una sintassi simile a JSP e PHP), \term{SquirrellyJS},
  \term{Nunjucks}, \term{LiquidJS}, \term{Eta} e \term{Haml} (per Ruby). Cambia
  la sintassi, non il concetto.}

\section{Analizzare il corpo delle richieste}

Gli utenti compilano il form della pagina per salvare nuovi elementi. Il form
viene inviato con il metodo \code{POST}, quindi dobbiamo analizzare il corpo
della richiesta per ottenere i nomi dei parametri e i loro valori. Non è così
semplice come sembra.

\begin{affianco}[0.58][c]
\begin{lstlisting}[style=pug, name=todo.pug]
form(action="/todo" method="POST")
  label(for="todo") Elemento da fare
  input(type="text" id="todo" name="todo" required)
  button(type="submit") Salva
\end{lstlisting}
\risultato
\begin{browser}[todo]
  Elemento da fare \wtcampo[2.6cm]{Learn Node.js}\ \wtbottone{Salva}
\end{browser}
\begin{lstlisting}[style=http, name=Richiesta generata]
POST /todo HTTP/1.1

todo=Learn+Node.js
\end{lstlisting}
\end{affianco}

\subsection{Il corpo arriva a pezzi}

Il corpo della richiesta \textbf{potrebbe non essere ancora disponibile} quando
iniziamo a elaborarla: magari l'utente ha una connessione lenta, o sta caricando
file di grandi dimensioni, e il corpo diventerà disponibile più tardi.

L'oggetto \code{request} (di tipo \code{http.IncomingMessage}) estende
\code{stream.Readable}: genera un evento \code{data} ogni volta che un nuovo
pezzo di dati diventa disponibile, e un evento \code{end} quando non c'è più
nulla da consumare. Per leggere il corpo dobbiamo quindi operare in modo
\textbf{asincrono}, sfruttando questi due eventi.

\begin{affianco}[0.55][c]
\begin{lstlisting}[style=js, name=body-parser.js]
function parseRequestBody(request) {
  return new Promise((resolve, reject) => {
    let body = [];
    request
      .on('data', chunk => { body.push(chunk); })
      .on('end', () => {
        body = Buffer.concat(body).toString();
        // qui `body` contiene l'intero corpo come stringa
        let data = parseRequestBodyString(body);
        resolve(data);
      });
  });
}
\end{lstlisting}
\risultato
\begin{immagine}[I pezzi arrivano con gli eventi \code{data};\\
                 all'evento \code{end} si ricompone il corpo.]
\begin{tikzpicture}[
    ch/.style={draw=primary!55!black, fill=primary!12!white, line width=0.6pt,
      font=\FiraCodeNL\scriptsize, minimum height=0.5cm, inner xsep=3pt},
    a/.style={-{Stealth[length=4pt,width=3.5pt]}, line width=0.7pt, draw=neutral!45!black}]
  \node[ch] (c1) at (0,0) {todo=Le};
  \node[ch] (c2) at (1.45,0) {arn+No};
  \node[ch] (c3) at (2.85,0) {de.js};
  \node[mcod, text=accent!70!black] (e) at (4.1,0) {end};
  \foreach \c in {c1,c2,c3} \node[mcod, anchor=south, text=accent!70!black, font=\FiraCodeNL\tiny] at (\c.north) {data};
  \node[ch, fill=green-gradient-1!45!white, draw=green-gradient-6] (b) at (2.0,-1.2)
    {"todo=Learn+Node.js"};
  \draw[a] (c2.south) -- node[mnota, right]{\code{Buffer.concat}} (b.north);
  \node[mnota, anchor=north] at (b.south) {poi \code{resolve(data)}};
\end{tikzpicture}
\end{immagine}
\end{affianco}

Si notino i due meccanismi che abbiamo studiato separatamente e che qui si
incontrano: gli \textbf{eventi} (capitolo 7) per raccogliere i pezzi, e una
\textbf{Promise} (capitolo 8) per esporre il risultato a chi ha chiamato.

\subsection{Interpretare la stringa}

Letto il corpo come stringa, si tratta di suddividerla per ottenere nomi e
valori dei parametri. Ricordiamo che il corpo è una stringa della forma
\code{p1=v1\&p2=v2\&...}, e che è \textbf{codificato come URL}.

\begin{affianco}[0.6][c]
\codicepiccolo
\begin{lstlisting}[style=js, name=body-parser.js]
function parseRequestBodyString(body) {
  let data = {};
  let slices = body.split("&");

  for (let slice of slices) {
    let paramName = slice.split("=")[0];
    paramName = decodeURIComponent(paramName.replace(/\+/g, " "));

    let paramValue = slice.split("=")[1];
    paramValue = decodeURIComponent(paramValue.replace(/\+/g, " "));

    data[paramName] = paramValue;
  }
  return data;
}
\end{lstlisting}
\risultato
\begin{immagine}[Dalla stringa codificata all'oggetto.]
\begin{tikzpicture}[
    st/.style={mcod, anchor=west},
    a/.style={-{Stealth[length=4pt,width=3.5pt]}, line width=0.7pt, draw=neutral!45!black},
    lab/.style={mnota, anchor=west, text=accent!70!black}]
  \node[st] (s0) at (0,0)     {"todo=Learn+Node.js\&when=oggi"};
  \node[st] (s1) at (0,-0.9)  {["todo=Learn+Node.js", "when=oggi"]};
  \node[st] (s2) at (0,-1.8)  {"todo" \quad "Learn Node.js"};
  \node[st] (s3) at (0,-2.7)  {\{ todo: "Learn Node.js", when: "oggi" \}};
  \draw[a] (0.3,-0.2) -- (0.3,-0.7) node[lab, midway, xshift=4pt]{split("\&")};
  \draw[a] (0.3,-1.1) -- (0.3,-1.6) node[lab, midway, xshift=4pt]{split("="), + \(\rightarrow\) spazio, decode};
  \draw[a] (0.3,-2.0) -- (0.3,-2.5) node[lab, midway, xshift=4pt]{data[nome] = valore};
\end{tikzpicture}
\end{immagine}
\end{affianco}

\warning{Ricordarsi della codifica}{
  Il corpo è codificato come URL, esattamente come descritto nel capitolo su
  HTML: la frase \guillemotleft Prove that P=NP\guillemotright\ viaggia come
  \code{Prove+that+P\%3DNP}. Bisogna quindi sostituire i \code{+} con spazi
  \textbf{e} applicare \code{decodeURIComponent}: saltare uno dei due passaggi
  produce bug che si manifestano solo con certi input.}

\section{Mettere insieme i pezzi}

Routing, template e lettura del corpo bastano per scrivere l'intera
applicazione. A lezione la si guarda dal vivo (il codice è fra i materiali del
corso); qui ne riportiamo una versione completa e funzionante, che si avvia con
\code{node app.js} dopo aver installato Pug. Gli elementi salvati stanno in un
array in memoria: si perdono a ogni riavvio del server, ma per ora va bene così.

Si comincia con lo stato dell'applicazione, il server e il router:

\begin{lstlisting}[style=js, name=app.js]
import http from 'http';
import fs from 'fs';
import pug from 'pug';

const PORT = 3000;
let todos = [];        // gli elementi salvati, in memoria
let nextId = 1;

let server = http.createServer();
server.listen(PORT);
server.on('request', handleRequest);

function handleRequest(request, response) {
  switch (request.url) {
    case "/":
      renderHomepage(request, response); break;
    case "/todo":
      handleTodoListRequest(request, response); break;
    case "/reset":
      handleResetRequest(request, response); break;
    case "/css/style.css":
      serveStatic(response, "./static/css/style.css", "text/css"); break;
    default:
      handleError(request, response, 404, "Pagina non trovata");
  }
}
\end{lstlisting}

I \textbf{controller}, uno per pagina, con la scelta in base al metodo per la
lista:

\begin{lstlisting}[style=js, name=app.js, firstnumber=last]
function renderHomepage(request, response) {
  render(response, "homepage");
}

function handleTodoListRequest(request, response, context = {}) {
  switch (request.method) {
    case "GET":
      handleTodoListRequestGet(request, response, context); break;
    case "POST":
      handleTodoListRequestPost(request, response, context); break;
    default:
      handleError(request, response, 405, "Metodo non supportato");
  }
}

function handleTodoListRequestGet(request, response, context) {
  render(response, "todo", { todos: todos, message: context.message });
}

async function handleTodoListRequestPost(request, response, context) {
  let data = await parseRequestBody(request);     // la funzione vista sopra
  let text = (data.todo || "").trim();
  if (text === "") {
    context.message = "Scrivi qualcosa prima di salvare.";
  } else {
    todos.push({ id: nextId++, text: text });
    context.message = "Elemento salvato!";
  }
  handleTodoListRequestGet(request, response, context);
}

function handleResetRequest(request, response) {
  todos = [];
  render(response, "reset");
}
\end{lstlisting}

E infine tre funzioni di servizio, per generare le pagine, servire i file
statici e rispondere con un errore:

\begin{lstlisting}[style=js, name=app.js, firstnumber=last]
function render(response, template, locals = {}, statusCode = 200) {
  let html = pug.renderFile(`./templates/${template}.pug`, locals);
  response.writeHead(statusCode, { "Content-Type": "text/html; charset=utf-8" });
  response.end(html);
}

function serveStatic(response, path, contentType) {
  fs.readFile(path, function (err, data) {
    if (err) {
      response.writeHead(500);
      return response.end();
    }
    response.writeHead(200, { "Content-Type": contentType });
    response.end(data);
  });
}

function handleError(request, response, statusCode, message) {
  render(response, "errorPage", { code: statusCode, description: message }, statusCode);
}
\end{lstlisting}

Le funzioni sono quelle viste una per una nelle sezioni precedenti; l'unica
novità è \code{render}, che raccoglie in un posto solo le tre righe che ogni
pagina ripeterebbe (generare l'HTML dal template, scrivere l'intestazione,
inviare il corpo). Il \code{context} passato da un gestore all'altro serve a
portare fino al template un messaggio per l'utente.

Tutte le pagine estendono lo stesso layout, che contiene la barra di
navigazione: è l'ereditarietà fra template vista prima, applicata a un caso
vero.

\begin{affianco}[0.5]
\begin{lstlisting}[style=pug, name=templates/layout.pug]
doctype html
html(lang="it")
  head
    meta(charset="utf-8")
    title To-do list
    link(rel="stylesheet" href="/css/style.css")
  body
    nav
      a(href="/") Home
      a(href="/todo") La lista
      a(href="/reset") Azzera
    main
      block content
\end{lstlisting}
\risultato
\begin{lstlisting}[style=pug, name=templates/todo.pug]
extends layout.pug

block content
  h1 La lista
  if message
    p.message #{message}
  form(action="/todo" method="POST")
    label(for="todo") Elemento da fare
    input(type="text" id="todo" name="todo" required)
    button(type="submit") Salva
  ul
    each todo in todos
      li #{todo.text}
    else
      li La lista è vuota.
\end{lstlisting}
\end{affianco}

\begin{esempio}{L'applicazione al lavoro}
Con il server avviato, si possono fare a mano le richieste che farebbe il
browser:
\begin{lstlisting}[style=shell, name=Terminale]
$ curl -X POST -d "todo=Learn+Node.js" http://localhost:3000/todo
...<p class="message">Elemento salvato!</p>...<ul><li>Learn Node.js</li></ul>...

$ curl -X POST --data-urlencode "todo=Prove that P=NP & <script>" http://localhost:3000/todo
...<ul><li>Learn Node.js</li><li>Prove that P=NP &amp; &lt;script&gt;</li></ul>...

$ curl -i -X DELETE http://localhost:3000/todo
HTTP/1.1 405 Method Not Allowed

$ curl -i http://localhost:3000/nonesiste
HTTP/1.1 404 Not Found
\end{lstlisting}
Il secondo elemento attraversa tutta la catena: \code{curl} lo codifica come
URL, \code{parseRequestBodyString} lo decodifica, e Pug, all'uscita, ne fa
l'escaping. Il \code{<script>} arriva al browser come testo, non come codice.
\end{esempio}

\Needspace{12\baselineskip}
\begin{infobox}{{Dopo una POST, ricaricare la pagina}}
  Dopo aver salvato un elemento, la pagina mostrata è la \textbf{risposta alla
  POST}. Se l'utente la ricarica, il browser chiede se reinviare il form: e
  l'elemento verrebbe salvato due volte. La soluzione abituale è il pattern
  \term{Post/Redirect/Get}: invece di mostrare la lista, il gestore della
  \code{POST} risponde con un reindirizzamento, e il browser chiede da solo la
  pagina con una \code{GET}, che si può ricaricare quante volte si vuole.
\begin{lstlisting}[style=js]
todos.push({ id: nextId++, text: text });
response.writeHead(303, { "Location": "/todo" });   // "vai a vedere /todo"
response.end();
\end{lstlisting}
\end{infobox}

\section{Estendere l'applicazione: cancellare un elemento}

Alla fine della lezione viene proposto un piccolo compito: aggiungere la
possibilità di \textbf{cancellare} un elemento, con gli strumenti visti finora.
I suggerimenti sono tre: aggiungere un nuovo percorso (per esempio
\code{/delete}), passare l'identificativo dell'elemento come \textbf{parametro
della query} (\code{/delete?id=X}), e mostrare accanto a ogni elemento della
lista un collegamento che punti all'URL giusto.

Prima di scrivere il gestore c'è un ostacolo. Il router confronta
\code{request.url} con stringhe esatte, ma \code{request.url} contiene
\textbf{anche la query}: per \code{/delete?id=3} vale proprio
\code{"/delete?id=3"}, che non è uguale a \code{"/delete"}, e la richiesta
finirebbe nel ramo \code{404}. Bisogna separare il percorso dai parametri, e
la classe \code{URL} lo fa per noi.

\begin{lstlisting}[style=js, name=app.js]
function handleRequest(request, response) {
  let url = new URL(request.url, `http://${request.headers.host}`);
  switch (url.pathname) {                       // solo il percorso, senza query
    /* ... le rotte di prima ... */
    case "/delete":
      handleDeleteRequest(request, response, url.searchParams); break;
    default:
      handleError(request, response, 404, "Pagina non trovata");
  }
}

function handleDeleteRequest(request, response, params) {
  let id = Number(params.get("id"));            // "3" -> 3
  let index = todos.findIndex(todo => todo.id === id);
  if (index === -1) {
    return handleError(request, response, 404, "Elemento inesistente");
  }
  todos.splice(index, 1);
  handleTodoListRequestGet(request, response, { message: "Elemento cancellato" });
}
\end{lstlisting}

Ora che il router guarda solo il percorso, anche \code{/todo?ordina=si} arriva
correttamente alla lista. Nel template ogni elemento riceve il suo
collegamento:

\begin{affianco}[0.55][c]
\begin{lstlisting}[style=pug, name=templates/todo.pug]
  ul
    each todo in todos
      li
        | #{todo.text}
        a(href=`/delete?id=${todo.id}`) cancella
    else
      li La lista è vuota.
\end{lstlisting}
\risultato
\begin{browser}[localhost:3000/todo]
  \textbullet\ Learn Node.js \wtlink{cancella}\\
  \textbullet\ Learn Pug \wtlink{cancella}\\
  \textbullet\ Learn Express \wtlink{cancella}
\end{browser}
\end{affianco}

Per questo usiamo un identificativo e non la posizione nell'array: dopo una
cancellazione le posizioni degli elementi successivi cambiano, e un collegamento
generato prima punterebbe all'elemento sbagliato. Un identificativo che non
esiste più produce invece un \code{404}, come è giusto.

\Needspace{14\baselineskip}
\begin{warningbox}{{Una GET non dovrebbe cambiare nulla}}
  La soluzione con un collegamento funziona, ma viola una regola di HTTP: una
  \code{GET} è un metodo \textbf{sicuro}, che non deve modificare lo stato del
  server (capitolo 9). Browser che precaricano i collegamenti, crawler dei motori
  di ricerca o un'estensione che controlla i link potrebbero «visitare»
  \code{/delete?id=3} e cancellare l'elemento senza che nessuno l'abbia chiesto.
  Lo stesso vale, a rigore, per la pagina \code{/reset}. La versione corretta usa
  un form con il metodo \code{POST}, e l'identificativo in un campo nascosto:
\begin{lstlisting}[style=pug]
li
  | #{todo.text}
  form(action="/delete" method="POST")
    input(type="hidden" name="id" value=todo.id)
    button(type="submit") cancella
\end{lstlisting}
  Il gestore legge allora \code{id} dal corpo con \code{parseRequestBody},
  esattamente come la \code{POST} su \code{/todo}, e risponde \code{405} a ogni
  altro metodo.
\end{warningbox}

\section{Anatomia di un'applicazione web}

Mettendo insieme i pezzi visti finora, si delinea l'architettura tipica di
un'applicazione web.

\begin{figure}[H]
  \centering
  \begin{tikzpicture}[font=\Lato\small,
      n/.style={wtnode, minimum width=3.0cm, minimum height=1.0cm},
      a/.style={-{Stealth[length=5pt,width=4pt]}, line width=0.85pt,
                draw=neutral!45!black}]

    \node[n, minimum height=1.8cm] (br) at (0,0)
      {\textbf{Web Browser}\\[3pt]\scriptsize HTML, CSS, JS};

    \node[n, draw=orange-gradient-6, fill=orange-gradient-1!40!white]
      (rt) at (4.6,1.3) {\textbf{Router}};
    \node[n] (bl) at (4.6,0) {\textbf{Logica di business}\\\scriptsize (controller)};
    \node[n] (pr) at (4.6,-1.4) {\textbf{Presentazione}\\\scriptsize (template)};
    \node[n, draw=green-gradient-6, fill=green-gradient-1!45!white]
      (db) at (8.8,0) {\textbf{Persistenza}\\\scriptsize dei dati};

    \draw[a] (br.east) to[out=25, in=180] node[wtlab, above, pos=0.5]
      {richiesta HTTP} (rt.west);
    \draw[a] (rt) -- (bl);
    \draw[a] (bl) -- (pr);
    \draw[a] (bl) -- (db);
    \draw[a, draw=accent!70!black] (pr.west) to[out=180, in=-25]
      node[wtlab, below, pos=0.5]{risposta HTTP} (br.east);

    \begin{scope}[on background layer]
      \node[rounded corners=5pt, draw=primary!40!white, line width=0.9pt,
            fit=(rt)(bl)(pr)(db), inner sep=11pt,
            label={[font=\Lato\scriptsize, text=primary!60!black]above:Backend}] {};
      \node[rounded corners=5pt, draw=primary!40!white, line width=0.9pt,
            fit=(br), inner sep=11pt,
            label={[font=\Lato\scriptsize, text=primary!60!black]above:Frontend}] {};
    \end{scope}
  \end{tikzpicture}
  \caption{L'architettura tipica di un'applicazione web tradizionale.}
  \label{fig:webapp-anatomy}
\end{figure}

\section{Riferimenti}

\begin{itemize}
  \item \textbf{Web Templating} --- EduTech wiki, Università di Ginevra. \\
        \urlx{https://edutechwiki.unige.ch/en/Web_templating}
  \item \textbf{Single page apps in depth}, Mikito Takada. \\
        \urlx{http://singlepageappbook.com/} \\
        Parte rilevante: \textit{Templating: from data to HTML}.
  \item \textbf{Node.js v.22.X (LTS) documentation}. \\
        \urlx{https://nodejs.org/docs/latest-v22.x/api/index.html} \\
        Per approfondire i metodi integrati di Node.js usati in questo capitolo.
  \item \textbf{A beginner's Guide to Pug}, James Hibbard. \\
        \urlx{https://www.sitepoint.com/a-beginners-guide-to-pug/}
\end{itemize}
